\documentclass[10pt,twocolumn]{article}
\usepackage[T1]{fontenc}
\usepackage[utf8]{inputenc}
\usepackage{lmodern}
\usepackage[margin=1.8cm,top=2.4cm,bottom=2cm,columnsep=0.6cm]{geometry}
\usepackage{fancyhdr}
\usepackage{titlesec}
\usepackage{booktabs}
\usepackage{amsmath,amssymb,amsfonts}
\usepackage{microtype}
\usepackage{xcolor}
\usepackage{mdframed}
\usepackage{parskip}
\usepackage{enumitem}
\usepackage{hyperref}

\definecolor{rulered}{RGB}{180,30,30}
\definecolor{boxgray}{RGB}{245,245,245}
\definecolor{keywordgray}{RGB}{100,100,100}

\pagestyle{fancy}
\fancyhf{}
\fancyhead[L]{\textsc{\small Claude Mag}}
\fancyhead[C]{\small\textit{Machine Learning Research Digest}}
\fancyhead[R]{\small 2026-09-02}
\fancyfoot[C]{\small\thepage}
\renewcommand{\headrulewidth}{0.8pt}

\titleformat{\section}{\normalsize\bfseries\color{rulered}}{}{0pt}{}%
  [\vspace{-4pt}\color{rulered}\rule{\columnwidth}{0.4pt}]
\titlespacing{\section}{0pt}{8pt}{4pt}

\hypersetup{colorlinks=true,urlcolor=rulered,linkcolor=black}

\begin{document}
\setlength{\parindent}{0pt}
\setlength{\parskip}{4pt}

{\small\color{keywordgray}\textbf{Keywords:}
  adversarial robustness \textbullet{}
  mixture-of-experts \textbullet{}
  mode connectivity \textbullet{}
  multi-norm defense}\par\vspace{4pt}

{\LARGE\bfseries\raggedright
  Robust CurveMoE: Multi-Norm Adversarial
  Defense for Mixture-of-Experts Models
  via Mode Connectivity}\par\vspace{4pt}

{\large\itshape\raggedright
  Low-Loss Bézier Paths Through MoE Weight Space
  Combine Per-Norm Robust Endpoints Into a
  Single Multi-Norm Robust Model}\par\vspace{6pt}

{\small\textbf{By} Xu Zhang, Ren Wang}\par\vspace{2pt}

{\small\color{keywordgray}arXiv:2608.26043 \textbullet{} Preprint, August 2026}
\par\vspace{2pt}\noindent{\color{rulered}\rule{\columnwidth}{0.6pt}}\vspace{6pt}

Mixture-of-Experts models are increasingly deployed at scale, but their
robustness to adversarial perturbations across multiple $\ell_p$ threat
models simultaneously has been understudied. This paper proposes Robust
CurveMoE: train separate MoE models robust under each threat norm, find
low-loss Bézier curve paths connecting them through weight space, and
sample a final model from the connecting path. The resulting model
achieves multi-norm adversarial robustness without ensemble inference
overhead, with clean accuracy dropping only 1.2\% relative to the
endpoint models.

\begin{mdframed}[backgroundcolor=boxgray,linecolor=rulered,linewidth=1pt,
    innerleftmargin=6pt,innerrightmargin=6pt,
    innertopmargin=6pt,innerbottommargin=6pt]
\textbf{\small AT A GLANCE}\par\vspace{4pt}
\begin{description}[leftmargin=0pt,labelsep=4pt,itemsep=2pt,parsep=0pt,
    font=\small\bfseries]
  \item[Problem:] MoE models trained for robustness under one $\ell_p$
    norm are vulnerable to other norms; union adversarial training is
    expensive and hurts accuracy.
  \item[Why it matters:] Production systems face diverse adversarial
    attacks; a single model must defend across all without multiplying
    inference cost.
  \item[Method:] Find Bézier curve paths through MoE weight space
    connecting per-norm robust endpoints; sample from the path at the
    point of best multi-norm trade-off.
  \item[Result:] Robust accuracy [86.2\%, 71.4\%, 68.3\%] under
    [$\ell_\infty$, $\ell_2$, $\ell_1$] vs.\ 4.8\% clean accuracy drop
    for union AT; CurveMoE drops only 1.2\%.
  \item[Evidence:] ResNet MoE on CIFAR-100; ablations on path degree
    and Gumbel-softmax relaxation.
\end{description}
\end{mdframed}

\section{The Big Picture}

Adversarial robustness is typically studied under a single threat model:
$\ell_\infty$-bounded perturbations with radius $\epsilon_\infty$. But
real-world adversaries can construct attacks in any $\ell_p$ geometry,
and a model robust to $\ell_\infty$ attacks may be easily fooled by
$\ell_2$ or $\ell_1$ perturbations.

Multi-norm robustness — simultaneous robustness to $\ell_\infty$,
$\ell_2$, and $\ell_1$ threats — has been addressed for dense networks
by union adversarial training or ensembles, but both approaches fail for
Mixture-of-Experts models:

\textbf{Union AT} trains on the worst-case perturbation over all three
norms simultaneously. This is computationally expensive and leads to
an accuracy-robustness trade-off that worsens as the number of norms
increases.

\textbf{Ensembles} combine predictions from three separately robust
models, tripling inference cost. For MoE models with thousands of
experts, ensemble inference is prohibitive.

Mode connectivity offers a third path: connect separately trained
robust models through low-loss paths in weight space and sample a model
from the path. For dense networks this has been demonstrated; for MoE
models, the routing-dependent loss landscape poses new challenges.

\section{Why Existing Approaches Fall Short}

\textbf{Single-norm AT:} Provides strong robustness under the target
norm but typically transfers poorly to orthogonal norms. A model
$\ell_\infty$-trained at $\epsilon = 8/255$ may have robust accuracy
below 40\% under $\ell_2$ attacks of matched strength.

\textbf{Union AT for MoE:} The multi-norm adversarial objective requires
solving $\max_{p \in \{\infty,2,1\}} \max_{\|\delta\|_p \leq \epsilon_p}
\mathcal{L}(\theta; x+\delta, y)$ at each training step. For MoE
models, the inner max interacts with discrete expert routing: the
worst-case perturbation may change the routing path, invalidating the
gradient from the current routing.

\textbf{Mode connectivity for dense networks:} Garipov et al.\ (2018)
showed that dense network pairs trained independently can be connected
by low-loss quadratic Bézier curves. For MoE models, the argmax
routing creates discontinuities that break the Garipov construction;
straight-line interpolation passes through high-loss regions.

\section{Core Mechanism}

The CurveMoE approach has four stages:

\begin{enumerate}[itemsep=2pt,parsep=0pt]
  \item \textbf{Per-norm endpoint training:} Train $\theta_{\ell_\infty}$,
    $\theta_{\ell_2}$, $\theta_{\ell_1}$ using standard adversarial
    training under each norm respectively.
  \item \textbf{Gumbel-softmax relaxation:} Replace argmax expert
    routing with Gumbel-softmax to make the loss differentiable with
    respect to weight interpolations along the connecting path.
  \item \textbf{Path optimization:} Find a Bézier curve $\phi(t)$
    connecting two endpoint models by minimizing the expected loss
    along the path under the relaxed routing.
  \item \textbf{Path sampling:} Evaluate a grid of $t$ values and
    select $t^*$ that maximizes the minimum robust accuracy across all
    three norms on a held-out validation set.
\end{enumerate}

For three endpoints, a triangular convex combination is used instead
of a pairwise curve, providing simultaneous connectivity to all three
robust models.

\section{Technical Formulation}

A degree-2 Bézier curve between $\theta_A$ and $\theta_B$ with
midpoint $\theta_c$ is:
\[
  \phi(t) = (1-t)^2 \theta_A + 2t(1-t) \theta_c + t^2 \theta_B,
  \quad t \in [0,1]
\]
The midpoint $\theta_c$ is optimized by minimizing:
\[
  \mathcal{L}_\text{path}(\theta_c) = \mathbb{E}_{t \sim U[0,1]}
    \sum_{p \in \{\infty,2,1\}} \alpha_p
    \mathcal{L}_p^\text{adv}(\phi(t))
\]
where $\mathcal{L}_p^\text{adv}(\theta) = \max_{\|\delta\|_p \leq \epsilon_p}
\mathcal{L}(\theta; x+\delta, y)$.

For the three-endpoint triangular path with $t_1 + t_2 \leq 1$:
\[
  \phi(t_1, t_2) = t_1 \theta_{\ell_\infty} + t_2 \theta_{\ell_2}
    + (1-t_1-t_2) \theta_{\ell_1}
\]
with three midpoint parameters $\theta_{c_{12}}, \theta_{c_{13}},
\theta_{c_{23}}$ along each edge of the triangle.

The Gumbel-softmax relaxation at temperature $\tau_g$:
\[
  \hat{g}_i = \frac{\exp((h_i + g_i)/\tau_g)}{\sum_j \exp((h_j + g_j)/\tau_g)}
\]
where $g_i \sim \text{Gumbel}(0,1)$ and $h_i$ are expert logits.
During path optimization $\tau_g = 1.0$; reverted to hard routing
($\tau_g \to 0$) for evaluation.

\section{Learning or Inference Procedure}

\begin{enumerate}[itemsep=2pt,parsep=0pt]
  \item Train $\theta_{\ell_p}$ for each $p \in \{\infty, 2, 1\}$
    using $\text{PGD}-10$ adversarial training under the respective norm.
  \item Initialize pairwise Bézier midpoints $\theta_c = (\theta_A + \theta_B)/2$.
  \item Optimize $\theta_c$ for 50 epochs with Gumbel-softmax routing
    at $\tau_g = 1.0$, Adam optimizer, learning rate $10^{-3}$.
  \item Grid-search $t \in \{0, 0.05, \ldots, 1.0\}$ on the validation set;
    select $t^*$ maximizing $\min_p R_p(\phi(t^*))$ where $R_p$ is robust
    accuracy under norm $p$.
  \item The final model is $\phi(t^*)$ with hard routing; no inference overhead.
\end{enumerate}

\section{What the Guarantee Says}

The path-loss objective guarantees that $\hat{V}(x_t)$ (the loss along
the optimized path) is uniformly low across $t \in [0,1]$. Because
the path connects models robust under each individual norm, and the
path loss is low (implying robustness is maintained along the path),
models at intermediate $t$ values inherit multi-norm robustness properties
from both endpoints.

This is not a formal robustness certificate but a constructive argument:
the path endpoint models have certified robustness under their respective
norms, and the path loss bounds how much that robustness degrades along
the path.

\section{Experimental Findings}

Evaluated on ResNet-56 MoE with 8 experts on CIFAR-100, under PGD-50
attacks with $\ell_\infty \epsilon = 8/255$, $\ell_2 \epsilon = 0.5$,
$\ell_1 \epsilon = 12$:

\begin{itemize}[itemsep=2pt,parsep=0pt]
  \item \textbf{CurveMoE robust accuracy:} [$\ell_\infty$: 86.2\%,
    $\ell_2$: 71.4\%, $\ell_1$: 68.3\%]
  \item \textbf{Best single-norm AT (per norm):} [$\ell_\infty$: 89.1\%,
    $\ell_2$: 74.2\%, $\ell_1$: 71.8\%] (different models for each)
  \item \textbf{Union AT:} [$\ell_\infty$: 81.4\%, $\ell_2$: 68.7\%,
    $\ell_1$: 65.2\%] with $-4.8\%$ clean accuracy drop
  \item \textbf{CurveMoE clean accuracy:} $-1.2\%$ relative to
    endpoint models
  \item \textbf{Path loss:} Maximum loss along the optimized path
    is 4\% higher than at endpoints; linear interpolation paths peak
    at 31\% higher, confirming the need for degree-2 curves.
\end{itemize}

\section{Ablations and Interpretation}

\textbf{Path degree:} Linear interpolation (degree 1) passes through
high-loss regions for MoE models. Degree-2 Bézier curves find a low-loss
path for all endpoint pairs tested. Degree 3 provides marginal improvement.

\textbf{Gumbel temperature:} $\tau_g = 1.0$ during path optimization
gives smooth gradients; lower temperatures ($\tau_g = 0.1$) approximate
hard routing but produce noisy gradients; $\tau_g = 0.5$ is a useful
intermediate for larger MoEs.

\textbf{Number of experts:} Path optimization scales well to 16 and
32 experts; beyond 64 experts, path optimization requires more epochs
to overcome routing instability.

\textbf{Transfer to ViT-MoE:} The approach transfers to Vision
Transformer MoE architectures with token routing; mode connectivity
persists, though the midpoint loss is 6\% higher than for ResNet MoE,
suggesting transformer MoE landscapes are less connected.

\section*{Reference}

Xu Zhang, Ren Wang.
\textbf{Robust CurveMoE: Multi-Norm Adversarial Defense for
Mixture-of-Experts Models via Mode Connectivity.}
arXiv:2608.26043, August 2026.
\url{https://arxiv.org/abs/2608.26043}

\end{document}
